\section{Discussion}

The central proposal of Paper 5 is that the trajectory of a research agent is itself a scientific object. A long-running agent does not merely output candidate theorems, code or literature summaries. It reveals which representations attract search, which memories are retrieved or missed, which failures recur, which local successes fail to compose, when verification changes the route, and which engineering constraints dominate the observed behaviour. If these traces are recorded as typed evidence rather than rewritten as narrative reflection, they can support explicit hypotheses about the research method.

This view changes the purpose of failure. In a conventional benchmark, a failed attempt is often discarded after scoring. In an open research programme, the same failure may be the most reusable result of the cycle. The live Millennium case study has already produced examples in which a route was narrowed because a surrogate failed to preserve the root coordinate, a source could not support the assumed transformation, a local theorem left the real interface open, or the research process itself missed relevant prior memory. None solves the root problem, but each can prevent future agents from spending identical effort under equivalent assumptions.

At the same time, storing failure is not sufficient. The strongest caution from both our architecture and recent memory research is that an abstraction can be worse than the episodes from which it was produced. Orion therefore makes lesson formation, failure diagnosis and tool promotion explicit transformations with contradiction, scope and validation obligations. The framework is designed to learn without requiring every learning event to become authoritative.

\subsection{Why metrology matters for claims about augmented intelligence}

The four-arm benchmark is perhaps the most important methodological addition for interpreting future results. ``An LLM with Orion solved the task'' is not evidence that Orion caused the solution. The model may already solve it; the framework may simply add tokens; a static checklist may explain the gain; or learned memory content may genuinely redirect search. Separating model-only, reset, sham-memory and learned conditions makes these explanations empirically distinguishable in a scoped task distribution.

This distinction is particularly important for scientific AI, where final success labels are sparse. A framework may contribute primarily by detecting false progress, reducing repeated failures or finding the missing interface earlier. These are legitimate effects, but they should be reported as epistemic-safety or efficiency effects rather than converted into a claim that the system became a stronger mathematician.

\subsection{Research architecture as an evolving experimental treatment}

We propose treating each method version as an intervention. Orion 3.1 should not be a marketing label for a larger repository; it should identify a challenger motivated by specific episodes, with a frozen mechanism and measurable prediction. This creates a sequence of natural experiments in which the method that generated the previous research record is itself changed and reevaluated.

The approach resembles automated agent design and self-modifying-agent work in maintaining variants and using empirical evaluation, but the scientific setting imposes additional constraints. The evaluator may be incomplete, truth and novelty are separate, fresh tasks are expensive, and a system capable of editing its own methods must not be allowed to redefine evidence after seeing outcomes. Protected assurance and governance therefore remain outside the candidate's authority.

\subsection{Practical implications of the case study}

For AI-system engineers, the project exposes a practical split between model capability and harness capability. Repository-native instructions, context compilation, retrieval coverage, state versioning, exact-subject CI, state migration and rollback are not peripheral details. They shape what the agent actually sees and what actions remain available. The live post-v3 CI repair is a useful example: a method-level architecture can be conceptually rich while repository integration still fails on mundane but load-bearing artifact and workflow details. Engineering incidents belong in the same longitudinal evidence system because they constrain reproducibility and deployment.

\subsection{What formal-methods and assurance researchers may ask}

The strongest assurance question is not whether every research decision can be proved optimal. It cannot. The more tractable question is whether authority escalation is structurally separated from generation. Paper 4 developed this idea for mathematical claims; Paper 5 applies the same philosophy to method evolution. A challenger may propose its own successor and may even discover an evaluator weakness, but a strong promotion claim depends on content-bound evidence and protected decision points outside that proposal path.

This architecture does not eliminate the trusted computing base. It makes it explicit. A compromised human administrator, CI service or governance credential can still override the system. The scientific claim is therefore scoped to observable process separation rather than absolute self-governance.

\subsection{What knowledge-representation researchers may ask}

The common substrate deliberately avoids treating all knowledge as one homogeneous graph. Epistemic claims, operators, episodes, obstructions and meta-method variants overlap in identity space while retaining specialized semantics and authority. Growth is measured on typed axes because a new relation and a new theorem are not interchangeable units. This may also make the longitudinal archive useful for studying whether reusable operator/strategy spaces saturate faster than the problem distribution.

The Orion-triviality hypothesis is intentionally empirical. If many verified tasks eventually require only stored, compositional or explicitly transferred structure, then increasing the atlas may shift effort from invention toward retrieval and gluing. If fresh tasks continue to demand new representations, operators or ontology, the hypothesis will be bounded or rejected. Either outcome informs the architecture.

\subsection{What scientists and mathematicians may ask}

The Millennium programmes are not presented as evidence that autonomous agents are close to solving the problems. Their role is to supply long-horizon, high-stakes research behaviour in domains where superficial progress is easy to overstate. Exact root contracts, primary-source discipline, typed failures, same-context-versus-independent review separation and explicit local-to-global obligations are intended to make this behaviour auditable.

A mathematically correct local lemma remains valuable even when it does not close the root. Conversely, an attractive representation with no faithful bridge to the root is treated as a research risk. The case study therefore produces a record not only of attempted solutions but of the methods by which a research system learns where not to place confidence.

\subsection{The empirical programme can fail productively}

Several plausible outcomes would force revision of our current design. Experience may cause fixation rather than improvement; lesson consolidation may harm more often than it helps; cross-domain memory may increase false transfer; the process overhead may exceed any saved search; or static Orion scaffolding may account for nearly all observed gains. A version may improve one domain while harming another. A coverage receipt may cost too much context. A root-coordinate audit may become excessively conservative. These are not loopholes in the paper; they are registered branches of the experiment.

The framework should therefore evolve only when the new evidence identifies a change whose benefit survives fresh testing. Recursive improvement without the possibility of a stable null result is not a scientific loop.

\section{Conclusion}

We introduced an evidence-governed formulation of continual research-method evolution. Orion represents immutable research episodes, derived lessons, failures, operators, motifs, problem fibres, gluing obstructions, saturation coordinates and framework variants as typed overlapping views of persistent external state. A live multi-agent mathematical programme provides an observational case study of the methods agents use, the ways they fail and the process defects they expose. The present evidence already motivates concrete framework hypotheses, but does not establish that a later version is superior.

To make that claim testable, Paper 5 contributes a governed upgrade protocol, repository-native coding-agent harness, typed process metrology, seven-axis lattice-growth measurement, four-arm causal-attribution benchmark, protected fresh assurance, operator-override semantics, variant archives and post-promotion attestation. Together these components define the experimental unit for future Orion evolution:
\rakldisp{
\boxed{
\text{observe}
\rightarrow\text{diagnose}
\rightarrow\text{propose}
\rightarrow\text{freeze}
\rightarrow\text{implement}
\rightarrow\text{compare}
\rightarrow\text{assure}
\rightarrow\text{govern}
\rightarrow\text{observe again}
}.
}

The next scientific result is therefore not predetermined. If future matched data show that Orion 3.1 or 3.2 reduces repeated research failures, improves fresh transfer, lowers false progress or reaches validated results more efficiently than the same underlying model and prior method under matched conditions, the evidence will support a scoped evolution claim. If they do not, the corresponding challenger should remain rejected or local. In either case, the longitudinal archive becomes an empirical record of how an AI-assisted research method changes, what it learns, and where its own theory of research fails.